\begin{titlingpage}
\centering
\vspace*{1.2in}
{\color{racing}\rule{\textwidth}{1.2pt}}\par\vspace{0.35in}
{\fontsize{30}{36}\selectfont\bfseries\color{racing} The Rhythm of Numbers\par}
\vspace{0.2in}
{\Large Continued Fractions, Primes, Music and Self-Similarity\par}
\vspace{0.35in}
{\color{racing}\rule{\textwidth}{1.2pt}}\par
\vspace{0.9in}
{\Large Teddy Rook\par}
\vspace{0.15in}
{\large Teddy Rook Book\par}
{\normalsize\url{https://teddyrookbook.com/connections/rhythm-of-numbers/}\par}
\vfill
{\large November 2026\par}
\vspace{0.4in}
\end{titlingpage}

\chapter*{Abstract}
\addcontentsline{toc}{chapter}{Abstract}

This monograph grows out of one guess: that $\pi$ is close to $3 + k/(7k+1)$ for a suitable whole number $k$. The guess is right at $k = 16$, where it gives $355/113$, the fraction Zu Chongzhi recorded in the fifth century AD, and it is right for a reason: $k/(7k+1)$ is a continued fraction cut off after two steps. We follow that reason in three directions.

First, \emph{rhythm}. Playing a number $\alpha$ as a beat (a hit wherever $\fl{n\alpha}$ steps up) produces a Sturmian word, and we give a self-contained proof that its blocks nest exactly as the continued fraction says: the blocks at level $n$ are $q_n$ beats long and hold $p_n$ hits. For $\pi - 3 = \cf{0; 7, 15, 1, 292, \dots}$ this gives bars of 7 beats, phrases of 106 and 113, and a single short phrase that begins at beat \ShortStart{} and closes a block of \BlockEnd{} beats. Reading the levels as bar, phrase, section and movement turns a continued fraction into musical form; Teddy Rook used it to compose the six-movement \emph{Continued-Fraction Suite} for piano, which we describe as a method.

Second, \emph{primes}. The same machine (the Euclidean algorithm) generates Euclidean rhythms, tunings from $\log_2(3/2)$, mode-locking staircases and the factoring method of Lehmer, Powers, Morrison and Brillhart. Coprimality ties these to $6/\pi^2$, and we derive the exact correlation of two pulse trains, $(\gcd(a,b)-1)/\sqrt{(a-1)(b-1)}$. We then pull two open threads with real computations. (a) \emph{Prime convergents.} Combining Lévy's theorem with the equidistribution of convergent denominators modulo $m$ (Szüsz; Moeckel; Jager and Liardet), we derive the heuristic that a typical real number has about $2\ln 2 \cdot \ln N \approx \CconstVal\ln N$ prime denominators among its first $N$ convergents. A sample of \RandM{} random reals fits the predicted growth; $\pi$ has \PiQPrimes{} prime denominators among its first $5{,}000$ convergents against \PiQExp{} expected, which is typical. (b) \emph{Spacings.} We place a crystal, the rhythm of $\pi$ (a one-dimensional quasicrystal with two gaps), \NZeros{} zeros of the Riemann zeta function, random GUE matrices, the primes and a Poisson process on one scale of rigidity, using nearest-neighbor spacings, pair correlation and number variance. The picture organizes the Montgomery--Dyson connection; it proves nothing about the Riemann hypothesis.

Third, \emph{self-similarity}. A continued fraction contains a copy of itself after every term, which is the idea behind Ford circles, the Stern--Brocot tree, Minkowski's question-mark function, the devil's staircase, continued-fraction Cantor sets (whose dimension we recompute as \DimEtwo{}), the Farey bulbs of the Mandelbrot set and the Fibonacci word fractal. A chi-square test of $\pi$'s first \NT{} terms against the Gauss--Kuzmin law gives $p = \PThirteen$ (12 degrees of freedom).

Returning to six open questions (Chapter~\ref{ch:openq}), we expand $\pi$ to twelve million certified terms, search its convergents for primes to $n = 20{,}000$, and prove that for a random real number the \emph{expected} number of prime denominators among the first $N$ convergents is $2\ln 2 \cdot \ln N + O(1)$, with exactly the heuristic's constant; the almost-everywhere version remains open. \addedon{\revA}

A second look (Chapter~\ref{ch:openq2}) closes that gap: for almost every real number the count of prime convergent denominators among the first $N$ convergents is $2\ln 2\cdot\ln N$ plus an error of order at most $(\ln N)^{1/2}(\ln\ln N)^{3/2+\varepsilon}$, so \cref{heur:primeconv} holds as a theorem for almost every number, though not, as far as anyone knows, for $\pi$. It also proves that $e$'s convergent denominators modulo a prime $\ell$ repeat with period dividing $6\ell$, shows that the top eigenvalue of the correlation matrix grows faster than any power of $\ln m$ and that its spectrum holds near-copies of small gcd matrices, explains the 171-note temperament by three small commas, and adds two tests of $\pi$'s terms. \addedon{\revB} A third look (Chapter~\ref{ch:third}), from Teddy Rook's notes, settles the growth of the correlation matrix's extreme eigenvalues: the logarithms of the largest eigenvalue and of the reciprocal of the smallest are both $(2 + o(1))\sqrt{\ln m/\ln\ln m}$. It proves that the eigenvalue distribution has a limit, and finds the limiting mean of $1/\lambda$. For $e$ it proves that the denominators change sign modulo an odd prime $\ell$ every $3\ell$ steps, so the period is exactly $6\ell$, and that a convergent denominator of $e$ can be prime only at $n = 2$ or $n \equiv 1 \pmod 3$. The search of $\pi$'s convergents, extended with a sieve that runs the recurrence modulo small primes to $n = 40{,}000$ for the denominators and $n = 28{,}000$ for the numerators, finds four new probable-prime denominators, at $n = 25{,}332$, $27{,}651$, $32{,}394$ and $37{,}721$, with \ThDigits{} to \FoDigits{} digits. They come in a cluster, and the prime convergents lean toward one note; both look like chance. Primes are far more plentiful among the intermediate fractions that belong to $\pi$'s largest terms. \addedon{\revC} Another chapter audits what is classical, what is known to specialists, what is a new presentation, and what is new. The mathematics is classical or known; what we did not find elsewhere is the reading of a number's continued fraction as a complete musical form, the Suite as an artwork, the prime-convergent experiment at this scale with an explicit constant, and the single spacing spectrum. Every number in the text was computed by the scripts listed in the appendices.

\vspace{1em}
\noindent\textbf{Keywords:} continued fractions; Sturmian words; Christoffel words; Euclidean rhythms; musical form; coprimality; prime numbers; Lévy's constant; pair correlation; random matrices; quasicrystals; fractals; Gauss--Kuzmin law.

\chapter*{How to read this work}
\addcontentsline{toc}{chapter}{How to read this work}

Each chapter is written to stand alone. Chapters~1--3 build the main tool. Chapters~4--7 follow it into music, coprimality and the primes, mostly retelling classical results with proofs or precise references. Chapters~8 and~9 are the two research threads, with new computations. Chapter~10 is about self-similarity and fractals, Chapter~11 about the music, Chapter~12 about the wider Connections diagram. Chapter~13 says what is new and what is not; Chapter~14 discusses limits and open questions, Chapter~15 reports further work on those questions \addedon{\revA}, Chapter~16 a second round of that work \addedon{\revB}, and Chapter~17 a third, with the digits of a new probable prime in Appendix~\ref{app:prime} \addedon{\revC}. The revision history below lists every addition.

Results that are classical carry a citation in their heading or the sentence that introduces them. Short standard proofs are included so that the text can be read by a strong undergraduate with calculus and a first course in number theory; longer ones are cited. A ``Heuristic'' is an argument that predicts a number without proving it, and we always say so. Interactive versions of most figures are on the companion pages \citep{RhythmSite,ConnectionsSite}.

Every number in this monograph was computed or checked by the scripts in Appendix~\ref{app:repro}. Numbers that come from data, from the terms of $\pi$ to the last decimal of a $p$-value, are inserted into the text automatically from the scripts' output; small numbers that appear in the prose (such as $355/113$ or $292 \times 113 = 32{,}996$) are asserted by the same scripts. All figures were drawn for this work from computed data.

\chapter*{Revision history}
\addcontentsline{toc}{chapter}{Revision history}

This work is updated as new results come in. Every addition carries the date it was published, in a small gray note such as \addedon{\revB}; text without a note belongs to the first version.

\begin{description}[style=nextline,leftmargin=2em]
\item[\revO{}] First version: Chapters 1--14 and Appendices A--C.
\item[\revA{}] Chapter~\ref{ch:openq}, ``Working on the open questions''; the status notes of Section~\ref{sec:open}; the abstract's paragraph on the open questions; six rows of the audit in Chapter~\ref{ch:new}; the long expansions in Appendix~\ref{app:methods} and six scripts in Appendix~\ref{app:repro}.
\item[\revB{}] Chapter~\ref{ch:openq2}, ``A second look at the open questions'': \cref{heur:primeconv} proved for almost every real number; the period of $e$'s denominators modulo a prime; growth bounds and near-copies in the spectrum of the correlation matrix; why the 171-note temperament wins; two further tests of $\pi$'s terms. Dated update notes in Section~\ref{sec:open}, in Chapter~\ref{ch:openq} and in the audit; one script in Appendix~\ref{app:repro}.
\item[\revC{}] Chapter~\ref{ch:third}, ``A third look at the open questions'', from Teddy Rook's notes: the growth of the largest and smallest eigenvalues of the correlation matrix; a limiting law for its spectrum; the sign flip and forced divisors of $e$'s convergents; $\pi$'s denominators searched to $n = 40{,}000$ and numerators to $n = 28{,}000$, with four new probable-prime denominators, the eleventh to fourteenth, whose digits are in Appendix~\ref{app:prime}; the cluster they form and the notes of the prime convergents; the rhythm sieve, prime intermediate fractions and the family of the opening guess. The chapter now follows Teddy Rook's finished text. A reader's map of what is proved (Section~\ref{sec:map}). Dated update notes in Chapters 5, 7, 8, 13, 14, 15 and 16, in the abstract and in the bibliography; new scripts in Appendix~\ref{app:repro}.
\end{description}

\cleardoublepage
\tableofcontents*
\clearpage
\listoffigures*
\clearpage
\listoftables*
